\uuid{6fT3}
\titre{Rayon de convergence et calcul de somme }

\niveau{L2}
\module{Séries entières}
\chapitre{Séries entières}
\sousChapitre{Rayon de convergence et calcul de somme}
\theme{Utilisation de la série exponentielle}
\auteur{Quentin Liard}
\datecreate{ 2026-04-10}
\organisation{AMSCC}
\difficulte{2}

\contenu{

\texte{ 
On considère la série entière à variable réelle $$S(x)=\displaystyle\sum_{n=0}^{+\infty}\frac{n+2}{n!}\,x^n.$$
}

\begin{enumerate}

\item   \question{Déterminer le rayon de convergence de la série entière $S(x)$ ainsi que son domaine de convergence.}
\reponse{
Posons
\[
a_n=\frac{n+2}{n!}.
\]
On étudie le quotient :
\[
\frac{a_{n+1}}{a_n}
=
\frac{n+3}{(n+1)!}\cdot \frac{n!}{n+2}
=
\frac{n+3}{(n+1)(n+2)}.
\]
Lorsque $n\to+\infty$,
\[
\frac{a_{n+1}}{a_n}\longrightarrow 0.
\]
Donc, pour tout $x\in\mathbb{R}$,
\[
\left|\frac{a_{n+1}x^{n+1}}{a_nx^n}\right|
=
|x|\left|\frac{a_{n+1}}{a_n}\right|
\longrightarrow 0.
\]
La série converge donc pour tout réel $x$.

Ainsi, le rayon de convergence est
\[
\boxed{R=+\infty}.
\]
L’intervalle ouvert de convergence est donc
\[
\boxed{]-\infty,+\infty[}.
\]
}

\item   \question{Déterminer la somme de la série sur l'intervalle ouvert de convergence.}
\reponse{
On écrit
\[
\sum_{n=0}^{+\infty}\frac{n+2}{n!}x^n
=
\sum_{n=0}^{+\infty}\frac{n}{n!}x^n
+
2\sum_{n=0}^{+\infty}\frac{x^n}{n!}.
\]
Or
\[
\sum_{n=0}^{+\infty}\frac{x^n}{n!}=e^x.
\]
De plus,
\[
\sum_{n=0}^{+\infty}\frac{n}{n!}x^n
=
\sum_{n=1}^{+\infty}\frac{x^n}{(n-1)!}.
\]
En posant $m=n-1$, on obtient
\[
\sum_{n=1}^{+\infty}\frac{x^n}{(n-1)!}
=
x\sum_{m=0}^{+\infty}\frac{x^m}{m!}
=
xe^x.
\]
Ainsi,
\[
\sum_{n=0}^{+\infty}\frac{n+2}{n!}x^n
=
xe^x+2e^x
=
(x+2)e^x.
\]

Donc la somme cherchée est
\[
\boxed{\sum_{n=0}^{+\infty}\frac{n+2}{n!}x^n=(x+2)e^x.}
\]
}

\end{enumerate}

}
